\practicechapter{第 14 章实践：C++ 与汇编互操作}

本章对应第一册第 14 章“C++ 与汇编互操作”。实践目标是理解 C++ 名字、符号名字、链接边界和反汇编之间的关系。你不需要一开始就写复杂汇编文件，但必须知道为什么 C++ 函数名到了二进制层可能变样，为什么 \code{extern "C"} 常被用作互操作边界。

\practicesection{一、对应原书问题：源码名字不是链接名字}

C++ 支持函数重载、命名空间、类成员函数等特性。为了让链接器区分这些函数，编译器会把源码名字编码成符号名字，这就是 \term{name mangling}{名字修饰}。例如两个同名但参数不同的函数在源码里都叫 \code{add}，链接层必须变成不同符号。

实践中，汇编、C、动态链接和工具脚本经常需要稳定符号名。此时可以用 \code{extern "C"} 关闭 C++ 名字修饰，让函数以 C ABI 风格暴露。它不是让函数变成 C 语言实现，而是改变链接名和调用边界约定。

\practicesection{二、从特例推到规律：同一个函数有多种名字}

特例一：源码中写 \code{int add(int,int)}，符号表里可能看到类似 \code{\_Z3addii} 的名字。特例二：用 \code{nm -C} 可以把修饰名 demangle 回更接近源码的形式。特例三：用 \code{extern "C" int add(int,int)} 后，符号表里更可能直接出现 \code{add}。

由此推出规律：报告中必须说明自己使用的是源码名字、mangled symbol name，还是 demangled display name。把三者混用，会导致你以为函数消失、重复或链接失败。

\practicesection{三、实践任务：比较 C++ 和 \code{extern "C"}}

写两个函数：

\begin{lstlisting}[language=C++]
int cpp_add(int a, int b) {
    return a + b;
}

extern "C" int c_add(int a, int b) {
    return a + b;
}
\end{lstlisting}

编译并观察：

\begin{lstlisting}[language=bash]
g++ -std=c++20 -O0 -g interop.cpp -o interop
nm ./interop > reports/interop.nm.txt
nm -C ./interop > reports/interop.nm-demangled.txt
readelf -s ./interop > reports/interop.symbols.txt
objdump -d ./interop > reports/interop.disasm.txt
\end{lstlisting}

再用 \code{cpu1ee\_cli} 输出 report，检查 symbol 数量上限是否影响观察。如果 report 中符号被 \code{max\_symbols} 截断，说明当前工具为了报告稳定性做了上限控制；需要用 readelf 或调整配置继续观察。

\practicesection{四、验收：互操作要讲清边界}

本章交付物是 \filepath{reports/ch14-cpp-assembly-interop.md}。通过标准如下：

\begin{enumerate}
  \item 能解释 name mangling 和 demangling 的区别。
  \item 至少比较一个普通 C++ 函数和一个 \code{extern "C"} 函数的符号名。
  \item 能说明符号名稳定为什么对汇编互操作有价值。
  \item 能指出 \code{extern "C"} 不等于关闭优化，也不等于改变函数体语言。
  \item 能把 \code{cpu1ee\_cli} 的 symbol 摘要和 \cmd{nm}/\cmd{readelf -s} 输出对应起来。
\end{enumerate}

这章的收获是：C++ 与汇编互操作不是“在 C++ 里夹一点汇编”这么简单，它首先要求你知道链接器看到的名字和边界是什么。
